\section{Experimental Setup}

\subsection{Evaluation Protocol}

All classical models use the same DS1 split (320 training and 80 test patients), with 30 predictors available before model-specific feature selection. The training classes are imbalanced (200 CKD and 120 non-CKD), so we use \glossterm{macrof1}{macro F1} as the main metric. We also report \glossterm{rocauc}{ROC AUC} and CKD \glossterm{recall}{recall}, since missing a CKD case is an important concern in screening. Hyperparameters for the Decision Tree, kNN and Naive Bayes were selected using 5-fold \emph{stratified} \glossterm{crossvalidation}{cross-validation} on the training set. We then compared different model versions using 10~$\times$~5-fold repeated stratified cross-validation and selected the version with the highest mean score. The test set was not used for model selection. Feature selection for kNN and Naive Bayes was performed within each fold. All models were implemented using scikit-learn\cite{pedregosa2011} with random seed 42.

\subsection{Classical Models}

\subsubsection{Decision Tree}\label{sec:dt}

We used a CART tree,\cite{breiman1984} which classifies patients through threshold-based splits on individual features. The fully grown tree fitted the training data perfectly (depth 5, 9 leaves), so we tested grid search and \glossterm{costcomplexitypruning}{cost-complexity pruning}. The grid covered \texttt{max\_depth} (1--8 or unlimited), \texttt{min\_samples\_leaf} (1, 2, 5, 10, 20), Gini or entropy\cite{quinlan1986}, and class weights, giving 180 combinations. The pruned tree ($\alpha = 0.0061$, depth 3, 4 leaves) achieved the highest repeated-CV score and was selected.

\subsubsection{k-Nearest Neighbours}\label{sec:knn}

kNN\cite{cover1967} predicts a patient's class from the nearest training patients, so the features need to be on a similar scale. We tested odd values of $k$ from 1 to 25, uniform and distance-weighted voting, and Manhattan and Euclidean distances. We also tested \glossterm{anovaftest}{ANOVA F-test} feature selection (5 to 30 features) and adding the hierarchical cluster label. The final model uses 15 features, $k=5$, Manhattan distance and distance-weighted voting.

\subsubsection{Naive Bayes}\label{sec:nb}

Gaussian Naive Bayes assumes that features are independent within each class. This assumption is difficult to meet in our dataset because some features are duplicated, such as creatinine and its log-transformed version. We tested different values of \texttt{var\_smoothing} ($10^{-12}$ to 1), ANOVA feature selection (1 to 30 features), and a mixed Gaussian--Bernoulli model for numerical and binary features. The final model uses three features (specific gravity, haemoglobin and packed cell volume) with \texttt{var\_smoothing} = 0.1.

\subsubsection{Random Forest}\label{sec:rf}

Random Forest\cite{breiman2001} combines 300 decision trees, each trained on a bootstrap sample with a random subset of features at each split. We used balanced class weights and the 15 features selected in Section~3, without further hyperparameter tuning. Its out-of-bag accuracy was 0.994. Feature importance was measured using mean decrease in impurity.

\subsection{Neural Networks}

\subsubsection{Convolutional Neural Network (CNN)}\label{sec:cnn}

We ran three initial CNN experiments on DS2, using a random stratified split of 70\% training, 15\% validation and 15\% test (1,389, 298 and 298 images). Images were resized with padding to 224 $\times$ 224 pixels, and class-weighted cross-entropy gave greater weight to the minority healthy class.

CNN V1 used three convolutional blocks followed by fully connected layers. V2 kept the same architecture but added image normalisation and random rotation (up to $\pm 5^\circ$). V3 added \glossterm{batchnormalisation}{batch normalisation} and adaptive average pooling, reducing the number of trainable parameters by about 97\% compared with V2.

All three models were trained for 20 epochs using Adam (learning rate 0.001) and a batch size of 32, without hyperparameter tuning. V1 was evaluated after training, while V2 and V3 used the checkpoints with the lowest validation loss. CNN V2 was selected for independent test evaluation because it achieved the best validation results.

\subsubsection{Later CNN Experiments (V4--V6)}\label{sec:tl}

Later experiments grouped consecutive IDs in blocks of 25 (80 groups; 1,399/300/286 training, validation and test images). Without patient or video IDs, this cannot confirm patient- or video-level separation.

V4 used ImageNet-pretrained ResNet18: its output layer was trained for two epochs (learning rate $10^{-3}$), then the full network for up to four ($10^{-4}$). V5 used our CNN with available bounding-box crops and stronger augmentation (flips, $\pm15^\circ$ rotation, brightness/contrast changes and blur); unboxed images received random training crops and centre crops at evaluation. V5 TL used pretrained ResNet18 (two head-only and up to six fine-tuning epochs). V6 added batch normalisation and global average pooling, reducing parameters from 12.9M to 94K; dropout was 0.3 and weight decay $10^{-4}$. All used balanced sampling, class-weighted cross-entropy (label smoothing 0.1), cosine scheduling and early stopping. Thresholds were tuned on validation data; test predictions averaged the original and horizontal/vertical flips.

